% !TeX root = ../main.tex
% appendices/bkv-migration.tex

\chapter{BKV Migration Reference}
\label{app:bkv-migration}

This appendix provides a component-by-component comparison between the
constructions described in Benton, Kennedy, and
Varming~\cite{BKV2009} and the corresponding implementations in our
library.  We compare against BKV's \emph{paper description}, not their
Coq library itself (which is not publicly available); where BKV show code, we
cite it, and where they describe a strategy in prose, we note what they
describe.

% ── Orders and CPOs ──────────────────────────────────────────────
\section*{Orders and CPOs}

\begin{table}[h]
\centering
\small
\begin{tabular}{@{}l l l@{}}
\toprule
\textbf{BKV description} & \textbf{Our implementation} & \textbf{Notes} \\
\midrule
\texttt{Record preorder}    & \texttt{Preorder.type} via \HB{}
  & No \texttt{.carrier} projections \\
CPO on preorder             & CPO requires \texttt{PartialOrder}
  & A\&J Def.\ 2.1.13 \\
\texttt{Class Pointed}      & \texttt{HasBottom} + \texttt{IsPointed}
  & \HB{} mixins \\
\texttt{lub\_of\_chain}     & \texttt{sup} ($\bigsqcup c$)
  & \\
\texttt{lub\_upper/least}   & \texttt{sup\_upper/sup\_least}
  & \\
Setoid equality \texttt{==} & Leibniz \texttt{=}
  & See DD-004 \\
\bottomrule
\end{tabular}
\end{table}

BKV build CPOs on preorders; we require antisymmetry
(\texttt{PartialOrder}), following Abramsky and
Jung~\cite{AJ1994}.  This is needed for \texttt{sup\_ext}
(equal chains have equal suprema).

% ── Continuous Functions ──────────────────────────────────────────
\section*{Continuous Functions}

\begin{center}
\small
\begin{tabular}{@{}l l l@{}}
\toprule
\textbf{BKV description} & \textbf{Our implementation} & \textbf{Notes} \\
\midrule
\texttt{fconti} record        & \texttt{cont\_fun} record & Renamed \\
\texttt{cf\_mfun} field       & \texttt{cf\_mono}         & Renamed \\
Separate \texttt{Continuous} module
  & \texttt{continuous} in \texttt{CPO.v}
  & Consolidated \\
\texttt{cont\_comp\_assoc}    & Same
  & via \texttt{proof\_irrelevance} \\
\bottomrule
\end{tabular}
\end{center}

We additionally define \texttt{strict\_fun} (continuous + bottom-preserving), a concept BKV do not discuss separately.

% ── Lift Monad ────────────────────────────────────────────────────
\section*{Lift Monad}

\begin{table}[h]
\centering
\small
\begin{tabular}{@{}l l l@{}}
\toprule
\textbf{BKV description} & \textbf{Our implementation} & \textbf{Notes} \\
\midrule
Coinductive \texttt{delay D}  & \texttt{option D}
  & Flat lift; see DD-006 \\
\texttt{Eps}/\texttt{Val}     & \texttt{None}/\texttt{Some}
  & \\
Coinductive order + bisim     & Flat order
  & \texttt{None} $\sqsubseteq$ all \\
Setoid equality (bisim)       & Leibniz \texttt{=}
  & \\
Sup described abstractly      & Proved via \texttt{ClassicalEpsilon}
  & 647 lines \\
\texttt{eta}/\texttt{bind}    & \texttt{ret}/\texttt{kleisli}
  & Direct pattern matching \\
\bottomrule
\end{tabular}
\end{table}

The coinductive \texttt{delay D} suffers from an antisymmetry
obstruction: \texttt{now d} $\neq$ \texttt{later (now d)} but they are
semantically equivalent under bisimulation.  Our supplementary file
\texttt{LiftMonad.v} develops this alternative fully and proves
the obstruction theorem; see \cref{sec:disc-lift-monad}.

% ── Function Spaces and Fixed Points ─────────────────────────────
\section*{Function Spaces and Fixed Points}

BKV describe the pointwise supremum of a chain of continuous functions
and note it is continuous, but the paper does not give a detailed
proof.
Our \texttt{FunctionSpaces.v} (882 lines) proves the full construction:
pointwise sup, its continuity, \texttt{curry}/\texttt{uncurry}/\texttt{eval},
and \texttt{FIXP} as a Scott-continuous operator.

For fixed points, BKV describe the iteration chain, the fixed-point
property, minimality, and Scott induction.  Our \texttt{FixedPoints.v}
(525 lines) proves all of these.

% ── Products and Sums ────────────────────────────────────────────
\section*{Products and Sums}

\textbf{Products.}  BKV describe binary products with pointwise order
and componentwise suprema.  Our \texttt{Products.v} (537 lines) gives
the full construction including continuous projections, pairing, and
the universal property.

\textbf{Sums.}  BKV describe separated sums but do not detail the
supremum construction.  Our \texttt{Sums.v} (624 lines) proves that
a chain in $A + B$ is eventually stable in one component, and the
supremum is the supremum of the stable projection.

% ── Domain Equations ──────────────────────────────────────────────
\section*{Domain Equations}

BKV describe solving $D \cong F(D, D)$ via inverse limits of
approximation sequences with EP-pairs.  Our
\texttt{DomainEquations.v} (446 lines) provides the abstract framework
with a single axiom (\texttt{bilimit\_exists}); all consequences
(ROLL/UNROLL, deflation chain) are fully derived.
\texttt{FunLift.v} (298 lines) registers the concrete bifunctor
$(A, B) \mapsto [A \to_c B]_\bot$.

% ── PCF Language ──────────────────────────────────────────────────
\section*{PCF Syntax and Semantics}

\begin{table}[h]
\centering
\small
\begin{tabular}{@{}l l l@{}}
\toprule
\textbf{BKV name} & \textbf{Our name} & \textbf{Notes} \\
\midrule
\texttt{TINT}/\texttt{TBOOL}   & \texttt{NLIT}/\texttt{BLIT}
  & Dropped \texttt{T} prefix \\
\texttt{TVAR i}                & \texttt{VAR v}
  & \\
\texttt{TFIX e}                & \texttt{FIX} $\tau_1\;\tau_2$ \texttt{body}
  & Explicit type indices \\
\texttt{TVAL}/\texttt{TLET}    & \texttt{VAL}/\texttt{LET}
  & \\
\texttt{TAPP}/\texttt{TFST}/\texttt{TSND}
  & \texttt{APP}/\texttt{FST}/\texttt{SND}
  & \\
\texttt{TOP}/\texttt{TGT}/\texttt{TIF}
  & \texttt{OP}/\texttt{GT}/\texttt{IFB}
  & \\
\bottomrule
\end{tabular}
\end{table}

We add the \texttt{var\_case} combinator for definitionally-computing
case analysis on de~Bruijn variables, avoiding the \texttt{JMeq\_eq}
opacity that BKV's \texttt{dependent destruction} introduces.

BKV describe the denotational semantics via point-free composition of
library combinators with setoid equality; we follow the same design
with Leibniz equality.  The substitution lemma is proved via an
explicit renaming chain (\texttt{sem\_ren\_ext} $\to$
\texttt{sem\_val\_ren}/\texttt{sem\_exp\_ren} $\to$
\texttt{sem\_val\_subst}/\texttt{sem\_exp\_subst}).

For adequacy, BKV describe the proof strategy (logical relation, Scott
induction in the FIX case) but characterise it as ``somewhat
involved.''
Our \texttt{PCF\_Adequacy.v} (820 lines) fully formalises this proof,
including admissibility for all four type cases and the fundamental
lemma by mutual induction.

% ── Axiom Comparison ─────────────────────────────────────────────
\section*{Axiom Usage}

BKV, being a paper, do not distinguish ``proved'' from ``axiomatic'' in
the Coq sense.  All our domain-theoretic suprema are fully proved:

\begin{table}[h]
\centering
\small
\begin{tabular}{@{}l l@{}}
\toprule
\textbf{Construction} & \textbf{Status} \\
\midrule
Function-space sup  & Proved constructively \\
Sum sup             & Proved constructively \\
Lift sup            & Proved via \texttt{ClassicalEpsilon} \\
Bilimit existence   & Axiom (\texttt{bilimit\_exists}) \\
\bottomrule
\end{tabular}
\end{table}

Classical axioms used: \texttt{ClassicalEpsilon} (in \texttt{Lift.v}),
\texttt{Classical} (in \texttt{PCF\_Adequacy.v} and
\texttt{ScottTopology.v}), \texttt{ProofIrrelevance} (in
\texttt{Morphisms.v}), and \texttt{FunctionalExtensionality} (in
\texttt{PCF\_Denotational.v} and \texttt{qCPOProperties.v}).
